\documentclass[11pt]{article}
\usepackage[a4paper,margin=0.82in]{geometry}
\usepackage{microtype}
\usepackage[T1]{fontenc}
\usepackage{lmodern}
\usepackage{parskip}
\usepackage[hidelinks]{hyperref}
\pagestyle{empty}
\setlength{\emergencystretch}{2em}
\setlength{\parskip}{0.45em}

\begin{document}

{\raggedright
\textbf{Cristiano Capone}\\
Istituto Superiore di Sanit\`a, 00161 Rome, Italy\\
\href{mailto:cristiano.capone@iss.it}{cristiano.capone@iss.it} \quad$\cdot$\quad \today\par}

\vspace{0.6em}
{\raggedright The Editors, \textit{Nature Communications}\par}

\vspace{0.8em}

Dear Editors,

We are pleased to submit our manuscript, ``Optimal stimulation sites are not the
most affected: personalised network models of Alzheimer's disease,'' by Cristiano
Capone, Enza Cece, Andrea Ciardiello, Guido Gigante,
Evaristo Cisbani and Maurizio Mattia, for consideration as an Article in
\textit{Nature Communications}.

Correcting altered brain-network dynamics can be attempted in two very different
ways, by reshaping the connectivity graph itself or by injecting focal activity
at a node, and these two interventions need not, and as we show do not, point to
the same target. In Alzheimer's disease, widely regarded as a distributed network
disorder, whether the connectivity signature can be reduced to a focal, stimulable
target has never been addressed causally, because doing so requires an
individualised generative model of brain dynamics that can be perturbed
\emph{in silico}. We fit subject-specific, cross-subject-identifiable
reservoir-computing models that reconstruct each individual's lagged functional
connectivity and use them as a testbed for such counterfactual interventions.

We make this distinction precise on the fitted model. Perturbing the connectivity
\emph{graph} (the fitted read-out kernel that defines the network's interactions)
and perturbing \emph{activity} (a focal oscillatory drive injected at a node)
prove to be fundamentally different interventions. Restoring control-like
dynamics by changing the graph is \emph{intrinsically distributed}: the required
correction is a coordinated, multi-site change of the connectivity kernel that
cannot be reproduced by editing any single node: a single-node graph edit is
inert. A focal activity perturbation, by contrast, \emph{can} be effective from a
single, physically realisable site.

Crucially, the site at which focal activity stimulation is optimal is \emph{not}
the site at which the graph is most altered. The ``graph'' sites (those with the
largest change in the connectivity kernel, predominantly weakly-connected
subcortical nodes) and the ``activity'' sites (those optimal for focal
stimulation, predominantly well-connected cortical nodes) are largely disjoint.
This dissociation is governed by the connectivity properties of the graph itself:
the most graph-perturbed nodes are the most weakly coupled, and therefore poor
actuators, whereas the nodes that provide therapeutic leverage are those best
positioned to redistribute a corrective perturbation network-wide. The optimal
target must therefore be chosen by its effect on the network's dynamics, not by
where the change in the graph concentrates. We further show that a real-time
closed-loop controller attains comparable efficacy at substantially lower dose
using only causally available information, the in-silico analogue of adaptive
neuromodulation.

These results reframe a central question in targeted neuromodulation, ``where
should we stimulate?'', as a question about the interaction between intervention
modality and network connectivity, answerable only with an individualised,
perturbable model. We believe this distinction between graph and activity
perturbation, grounded in the graph's connectivity properties, will interest the
broad readership of \textit{Nature Communications} across network neuroscience,
computational modelling and neuromodulation, well beyond Alzheimer's disease.

The manuscript is original, is not under consideration elsewhere, and all authors
have approved its submission. The data derive from the Alzheimer's Disease
Neuroimaging Initiative and are used under its Data Use Agreement. We declare no
competing interests, and would gladly suggest suitable referees on request. We
thank you for considering our work.

\vspace{1em}
{\raggedright
Sincerely,\\[0.3em]
Cristiano Capone, on behalf of all authors\par}

\end{document}
